% Generated by scripts/generate-cv.mjs from src/data/profile.json and publications.json.
% Standalone, Overleaf-compatible. Upload this file; compile with pdfLaTeX.
% Edit canonical data and regenerate instead of editing this generated file.
\documentclass[11pt,letterpaper]{article}
\usepackage[margin=0.72in]{geometry}
\usepackage[T1]{fontenc}
\usepackage[utf8]{inputenc}
\usepackage{lmodern}
\usepackage{microtype}
\usepackage{tabularx}
\usepackage{enumitem}
\usepackage{titlesec}
\usepackage{fancyhdr}
\usepackage{xcolor}
\usepackage{hyperref}
\ifdefined\pdfgentounicode
  \input{glyphtounicode}
  \pdfgentounicode=1
\fi
\definecolor{linkblue}{HTML}{244E7A}
\hypersetup{colorlinks=true,urlcolor=linkblue,linkcolor=linkblue,
  pdftitle={Kaiming Liu -- Curriculum Vitae},pdfauthor={Kaiming Liu}}
\urlstyle{same}
\setlength{\parindent}{0pt}
\setlength{\parskip}{3pt}
\setlength{\emergencystretch}{2em}
\renewcommand{\baselinestretch}{1.03}
\titleformat{\section}{\large\bfseries}{}{0pt}{}
\titlespacing*{\section}{0pt}{13pt}{6pt}
\setlist[itemize]{leftmargin=15pt,itemsep=3pt,topsep=4pt,parsep=0pt}
\pagestyle{fancy}
\fancyhf{}
\renewcommand{\headrulewidth}{0pt}
\fancyfoot[L]{\footnotesize Kaiming Liu \quad | \quad Updated October 1, 2026}
\fancyfoot[R]{\footnotesize \thepage}
\setlength{\footskip}{24pt}
\newcommand{\entry}[3]{%
  \noindent\begin{tabularx}{\linewidth}{@{}Xr@{}}\textbf{#1}&{\small #3}\end{tabularx}\par\nobreak
  {\itshape #2}\par\nobreak}
\begin{document}
{\LARGE Kaiming Liu}\par
Physics PhD student, University of Illinois Urbana-Champaign\par
\href{mailto:kl54@illinois.edu}{kl54@illinois.edu} \enspace $\cdot$ \enspace
\href{https://kl543.github.io/}{kl543.github.io} \enspace $\cdot$ \enspace
\href{https://github.com/kl543}{GitHub} \enspace $\cdot$ \enspace
\href{https://www.linkedin.com/in/kaiming-l-9a166526a/}{LinkedIn}

\section*{Research Interests}
Scientific machine learning; inverse problems; probabilistic and multimodal inference; physics-based forward validation; simulation-to-experiment transfer; computational condensed matter physics.

\section*{Education}
\entry{University of Illinois Urbana-Champaign}{Ph.D. in Physics, in progress}{Aug 2023--present}
Advisor: Bryan K. Clark\par
\vspace{4pt}
\entry{Xi'an Jiaotong University}{B.S. in Physics (Honors)}{Sep 2019--Jul 2023}

\vspace{4pt}
\entry{University of California, Berkeley}{Berkeley Physics International Education (BPIE) Program}{Aug--Dec 2021}
Non-degree visiting program\par
\vspace{4pt}
\entry{Xi'an Jiaotong University}{Special Class for the Gifted Young}{Sep 2017--Jul 2019}

\vspace{4pt}

\section*{Research Experience}
\entry{University of Illinois Urbana-Champaign}{Graduate research, Bryan K. Clark group}{Nov 2025--present}
\begin{itemize}
  \item Develop probabilistic approaches to crystal electric field inverse problems, including multimodal parameter inference and reconstruction of observables with physics-based forward models.
  \item Study the sensitivity of scientific inference to experimental noise and simulation-to-experiment distribution shift, with applications to magnetic and scattering models.
  \item Implemented neural-network backflow wavefunctions for strongly correlated fermionic systems and ran variational Monte Carlo simulations with parallel sampling on computing clusters.
  \item Profiled sampling, local-energy evaluation, and optimization stages; refactored sampling routines to reuse wavefunction evaluations.
\end{itemize}
\vspace{3pt}
\entry{University of Illinois Urbana-Champaign}{Research assistant, Vidya Madhavan group}{Aug 2023--Oct 2025}
\begin{itemize}
  \item Grew thin-film heterostructures by molecular beam epitaxy and characterized materials using STM/STS in ultrahigh-vacuum systems.
  \item Built an automated MBE temperature-control program with feedback loops, measurement logging, and alerts to support repeatable growth protocols.
  \item Developed Python and MATLAB analysis routines for STM/STS data, including constant-energy contour extraction and joint-density-of-states calculations for quasiparticle-interference studies.
\end{itemize}
\vspace{3pt}
\entry{Xi'an Jiaotong University}{Research assistant, Yongchang Zhang group}{Mar 2022--Jul 2023}
\begin{itemize}
  \item Simulated Rydberg-atom dynamics and photon-mediated interactions with MATLAB and Python to study quantum nonlinear optics.
  \item Built a discrete-time lattice simulation to examine two-qubit phase-flip dynamics and explore many-body regimes.
\end{itemize}
\vspace{3pt}
\entry{University of California, Berkeley}{Research assistant, Michael F. Crommie group}{Oct--Dec 2021}
\begin{itemize}
  \item Assisted with synthesis and STM characterization of graphene nanoribbons containing five-membered rings, supporting studies of their electronic states.
\end{itemize}
\vspace{3pt}

\section*{Publications}
Zhen Zhu, Yudi Huang, Julian May-Mann, \textbf{Kaiming Liu}, Zheyu Wu, Shanta R. Saha, Johnpierre Paglione, Alexander G. Eaton, Andrej Cabala, Michal Vali\v{s}ka, Eduardo Fradkin, Vidya Madhavan.\par
\href{https://doi.org/10.1073/pnas.2602117123}{Evidence of intertwined pair density and charge density wave orders in UTe2}. \textit{PNAS} 123 (29), e2602117123 (2026). \href{https://arxiv.org/abs/2603.08688}{arXiv version}.\par



\section*{Teaching}
\entry{University of Illinois Urbana-Champaign}{Teaching assistant, PHYS 102}{Aug--Dec 2023}
Led discussion and laboratory sections, held office hours, and provided feedback on assignments in electromagnetism and modern physics.\par

\section*{Selected Honors}
\textbf{China National Scholarship} \hfill 2021\\
Xi'an Jiaotong University\par
\textbf{Mathematical Contest in Modeling, Meritorious Winner} \hfill 2021\\
COMAP\par

\section*{Technical Skills}
\textbf{Programming and computing:} Python, NumPy, SciPy, JAX, MATLAB; Git; Linux/bash; parallel cluster computation.\par
\textbf{Scientific methods:} Probabilistic inference, Monte Carlo (MCMC/VMC), numerical optimization, forward-model validation, runtime profiling.\par
\textbf{Experiment:} Molecular beam epitaxy, STM/STS, ultrahigh-vacuum systems, feedback control, scientific data analysis.\par
\end{document}
